\chapter{Dynamic Classic Algorithms and Core Operators}

本章从静态结构向动态算子演进，涵盖设计范式、排序、查找与图论算法。数学层只谈规约与定理，\texttt{Visit}、优先队列、并查集实现均为实例化。

\section{顶层设计范式}

\subsection{分治策略 (Divide and Conquer)}
\begin{itemize}
    \item \textbf{对象}：规模 $n$ 的问题 $P$ 划分为 $a$ 个规模 $n/b$ 的子问题。
    \item \textbf{递推}：$T(n)=aT(n/b)+f(n)$，由主定理决定渐进阶。
    \item \textbf{实例化}：归并、快排为分治实例。
\end{itemize}

\subsection{贪心算法 (Greedy)}
每步取局部最优 $s_i=\arg\max_{x\in C}\mathcal{F}(x)$。全局最优需：
\begin{enumerate}
    \item \textbf{最优子结构}：$P$ 的最优解含子问题 $P'$ 的最优解。
    \item \textbf{贪心选择性}：局部最优序列可拼接为全局最优。
\end{enumerate}

\newpage

\section{排序算法算子 (Sort)}

输入 $A=\langle a_1,\dots,a_n\rangle$，目标 $A'=\langle a'_1,\dots,a'_n\rangle$ 使 $a'_1\le\cdots\le a'_n$。

\subsection{比较类排序}
\begin{enumerate}
    \item \textbf{冒泡排序}：相邻交换消逆序对 $I(A)=|\{(i,j)\mid i<j\land a_i>a_j\}|$。
    \item \textbf{选择排序}：$a'_k=\min\{a_k,\dots,a_n\}$，比较次数固定 $n(n-1)/2$。
    \item \textbf{插入排序}：维护已序前缀 $A[1..k]$，插入 $a_{k+1}$；已序时 $O(n)$。
    \item \textbf{快速排序}：基准 $p$ 划分 $A_L=\{x\le p\},A_R=\{x>p\}$，递归。
    \item \textbf{归并排序}：$T(n)=2T(n/2)+O(n)$，稳定 $O(n\log n)$，需 $O(n)$ 辅助。
    \item \textbf{堆排序}：建堆 $O(n)$，反复取根下滤，$O(n\log n)$，$O(1)$ 辅助。
\end{enumerate}

\begin{algorithm}[H]
\caption{快速排序规约}
\begin{algorithmic}[1]
\Procedure{QuickSort}{$A,l,r$}
    \If{$l\ge r$} \Return \EndIf
    \State $p\gets\text{Partition}(A,l,r)$
    \State \text{QuickSort}$(A,l,p-1)$; \text{QuickSort}$(A,p+1,r)$
\EndProcedure
\Procedure{Partition}{$A,l,r$}
    \State $pivot\gets A[r]$; $i\gets l-1$
    \For{$j\gets l$ \textbf{to} $r-1$}
        \If{$A[j]\le pivot$} $i\gets i+1$; \textbf{Swap}($A[i],A[j]$) \EndIf
    \EndFor
    \State \textbf{Swap}($A[i+1],A[r]$); \Return $i+1$
\EndProcedure
\end{algorithmic}
\end{algorithm}

\subsection{非比较类排序}
\begin{enumerate}
    \item \textbf{计数排序}：$C[x]=|\{a_i\mid a_i=x\}|$，前缀 $S[x]=\sum_{k\le x}C[k]$ 定位置，$O(n+k)$。
    \item \textbf{桶排序}：值域 $[Min,Max]$ 分 $k$ 区间 $B_1\dots B_k$，桶内排序后串联。
    \item \textbf{基数排序}：按位 $d=0\dots D-1$ 稳定排序（以计数为子程序）。
\end{enumerate}

\begin{itemize}
    \item \textbf{实例化}：网页 \texttt{array-algorithms} 与 \texttt{race} 分别演示各排序的帧与实测对比；\texttt{Swap} 为语言实例。
\end{itemize}

\newpage

\section{查找与并查集}

\subsection{二分查找 (Binary Search)}
\begin{itemize}
    \item \textbf{对象}：有序数组 $A$ 上区间 $[l,r]$，$m=\lfloor(l+r)/2\rfloor$。
    \item \textbf{不变式}：$A[m]=x\implies$找到；$A[m]<x\implies l=m+1$；否则 $r=m-1$。
    \item \textbf{定理}：时间 $O(\log n)$，需随机访问 $O(1)$（顺序存储实例）。
    \item \textbf{实例化}：网页 \texttt{array-algorithms:search} 演示区间收缩。
\end{itemize}

\subsection{并查集 (Union-Find)}
\begin{itemize}
    \item \textbf{对象}：不相交集合族 $\mathcal{S}=\{S_1,\dots,S_k\}$。
    \item \textbf{操作}：$\text{Find}(x)$ 取代表元，$\text{Union}(x,y)$ 合并；按秩合并+路径压缩后摊还 $O(\alpha(n))$。
    \item \textbf{实例化}：图论 Kruskal、Tarjan 中复用本结构；网页 \texttt{graph-mst} 演示并查集判环。
\end{itemize}

\newpage

\section{图论核心算法}

\subsection{最短路径算法}
维护 $D:V\to\mathbb{R}\cup\{\infty\}$，松弛：$D[u]+w(u,v)<D[v]\implies D[v]\gets D[u]+w(u,v)$。

\begin{enumerate}
    \item \textbf{Dijkstra}：贪心取 $u\in V\setminus S$ 使 $D[u]$ 最小入 $S$ 并松弛。不支持负权；优先队列 $\text{ExtractMin}/\text{DecreaseKey}$ 为贪心选择的实例化（实例取堆，见树章堆序）。
    \item \textbf{Bellman-Ford}：$|V|-1$ 轮全量松弛；第 $|V|$ 轮可松弛 $\implies$ 负环，$O(|V||E|)$。
    \item \textbf{Floyd-Warshall}：$D^{(k)}[i][j]=\min(D^{(k-1)}[i][j],D^{(k-1)}[i][k]+D^{(k-1)}[k][j])$，$O(|V|^3)$，支持负权，$D[i][i]<0\implies$负环。
\end{enumerate}

\begin{algorithm}[H]
\caption{Dijkstra 规约（$Q$ 为优先队列实例）}
\begin{algorithmic}[1]
\State \textbf{输入}: $G=(V,E),s,w$; $\forall v,D[v]\gets\infty,D[s]\gets0$; $S\gets\emptyset$; $Q\gets V$
\While{$Q\neq\emptyset$}
    \State $u\gets\text{ExtractMin}(Q)$; $S\gets S\cup\{u\}$
    \ForAll{$v\in N(u)$} \If{$D[u]+w(u,v)<D[v]$} $D[v]\gets D[u]+w(u,v)$; $\text{DecreaseKey}(Q,v,D[v])$ \EndIf \EndFor
\EndWhile
\end{algorithmic}
\end{algorithm}

\begin{itemize}
    \item \textbf{实例化}：网页 \texttt{graph-dijkstra}、\texttt{graph-bellman}、\texttt{graph-floyd} 分别演示三算法的 $D$ 表；堆为 Dijkstra 的 PQ 实例。
\end{itemize}

\subsection{最小生成树 (MST)}
\begin{enumerate}
    \item \textbf{Prim}：任意 $r\in V$，维护 $T$，取横切边最小者入 $T$（割定理），$O(|E|\log|V|)$。
    \item \textbf{Kruskal}：$E$ 升序，依序取 $(u,v)$ 若 $\text{Find}(u)\neq\text{Find}(v)$ 则入 MST（并查集实例），$O(|E|\log|E|)$。
\end{enumerate}

\subsection{拓扑排序：Kahn 算法}
\begin{itemize}
    \item \textbf{对象}：DAG 的入度 $\text{indeg}(v)=|\{u\mid(u,v)\in E\}|$。
    \item \textbf{规约}：$Q=\{v\mid\text{indeg}(v)=0\}$；弹出 $u$ 入 $\sigma$，$\forall v\in N(u),\text{indeg}(v)--$ 为 $0$ 则入 $Q$；$|\sigma|<|V|\implies$有环。
    \item \textbf{实例化}：$Q$ 的公平性为实例（队列/任意集均正确）；网页 \texttt{graph-toposort} 演示入度表。
\end{itemize}

\subsection{启发式搜索：A* (A-Star)}
\begin{itemize}
    \item \textbf{对象}：加权图 $G$，起点 $s$，终点 $t$，启发式 $h:V\to\mathbb{R}_{\ge0}$ 估计 $v\to t$ 代价。
    \item \textbf{不变式}：$g(v)$ 为 $s\to v$ 已知最短，$f(v)=g(v)+h(v)$ 为优先键；\textbf{可采纳} $h(v)\le h^*(v)$（真实最优），\textbf{一致} $h(u)\le w(u,v)+h(v)$。
    \item \textbf{定理}：可采纳时 A* 首次出队 $t$ 即最优；$h\equiv0$ 退化为 Dijkstra；一致时无需重开集。
    \item \textbf{实例化}：网页 \texttt{graph-astar} 以 $f$ 标注演示 $h=0$/曼哈顿/欧几里得模拟；真实坐标 $h$ 为实例化。
\end{itemize}

\subsection{树上查询：LCA 倍增}
\begin{itemize}
    \item \textbf{对象}：有根树 $T$，$depth[v]$，$up[v][j]=v$ 的 $2^j$ 祖先，$LOG=\lceil\log_2|V|\rceil$。
    \item \textbf{不变式}：$up[v][0]=parent(v)$，$up[v][j]=up[up[v][j-1]][j-1]$。
    \item \textbf{规约}：查询 $(u,v)$：(1) 提深者至同深；(2) 同步上跳至分叉点之子，父即 LCA。预处理 $O(|V|\log|V|)$，查询 $O(\log|V|)$。
    \item \textbf{实例化}：网页 \texttt{graph-lca} 演示深度表与二进制提升路径。
\end{itemize}

\section{高级工程专项算子}

\subsection{网络流 (Network Flow)}
\begin{itemize}
    \item \textbf{对象}：$G=(V,E)$，容量 $c:E\to\mathbb{R}^+$，流 $f:E\to\mathbb{R}$ 满足 $0\le f\le c$ 且 $\forall u\neq s,t,\sum_v f(u,v)=0$。
    \item \textbf{残量}：$c_f(u,v)=c(u,v)-f(u,v)$，增广路 $P$ 在 $G_f$ 中 $s\to t$。
    \item \textbf{定理（最大流最小割）}：$|f|_{\max}=\min_{S\ni s,T\ni t} c(S,T)$；增广至无 $s\to t$ 路即最优。
    \item \textbf{实例化}：\textbf{Edmonds-Karp} 以 BFS 找增广 $O(|V||E|^2)$；\textbf{Dinic} 以 BFS 分层 + DFS 阻塞流 $O(|V|^2|E|)$（二分图 $O(\sqrt{|V|}|E|)$），为同一残量定理的不同实例；网页 \texttt{graph-maxflow} 演示 Dinic 分层与阻塞流。
\end{itemize}

\subsection{二分图最大匹配}
\begin{itemize}
    \item \textbf{对象}：$G=(X\cup Y,E)$，匹配 $M\subseteq E$ 无公共点。
    \item \textbf{定理（Berge）}：$M$ 最大 $\iff$ 无增广路 $P$（端点未匹配、交替）。$M'\gets M\oplus P$ 增大。
    \item \textbf{实例化}：匈牙利算法反复找增广路；网页虽未单列模块，但为流的特例（建源汇容量 $1$）。
\end{itemize}

\subsection{强连通分量 (SCC)：Tarjan 算法}
\begin{itemize}
    \item \textbf{对象}：DFS 时间戳 $DFN[u]$，追溯 $LOW[u]$。
    \item \textbf{不变式}：$u$ 入栈；$v$ 未访则 $LOW[u]=\min(LOW[u],LOW[v])$；$v$ 在栈则 $LOW[u]=\min(LOW[u],DFN[v])$。
    \item \textbf{定理}：$LOW[u]=DFN[u]\iff u$ 为 SCC 根，栈中 $u$ 及以上一次弹出构成一 SCC，$O(|V|+|E|)$。
    \item \textbf{实例化}：网页 \texttt{graph-scc} 演示 $DFN/LOW$ 与弹栈帧。
\end{itemize}
